\section{Perception and Temporal Obstacle Estimation}
\label{sec:estimation}

\subsection{Common Interface and Qualification Layer}

Comparing temporal estimators fairly requires that they be interchangeable
without changing anything else, so all variants live behind one interface and
inside one node. A detection source publishes raw image-space obstacles; the
temporal estimator consumes them and republishes a planner-facing stream on
which bearing, apparent area, and risk are recomputed by a single shared
formula; the planner is unchanged across variants. The passthrough variant T0
runs inside the same node and launch as the filtered variants, so runtime
differences cannot bias the comparison.

Three upstream conditions must be distinguished, and conflating two of them
caused a documented closed-loop failure before the distinction was made
explicit. A \emph{fresh detection} is new information and updates the state.
A \emph{fresh empty message} is evidence of absence: the detector ran and saw
nothing, which increments a miss counter and eventually deletes the track.
\emph{Silence} is the absence of information: no message arrived at all,
which is bridged by prediction up to a bounded horizon. Treating silence as
absence caused a planner to abandon a committed maneuver and recover onto a
wrongly recaptured route; treating absence as silence produces ghost
obstacles. Both failure modes are avoided by handling the three cases
separately, and the distinction is tested at unit and replay level.

A common qualification layer precedes all estimator variants, so that no
variant can benefit from unqualified detections. It admits a track only after
a warm-up of 20 stream-healthy updates spanning at least 1.0\,s, applies
physics-derived coherence bounds (image-centre rate $\leq 1.2$\,units/s,
$|\Delta \log \text{size}| \leq 0.5$ per message), maintains a reference with
a 1.0\,s maximum age reset after five consecutive rejections, and confirms a
track only after $M=3$ accepted coherent updates within a 1.5\,s window. The
source-restart silence threshold is 4.0\,s, deliberately longer than the
2.5\,s estimator horizon so that a bridged dropout can never trigger a
mid-maneuver re-warm-up. These parameters were frozen before the reported
campaign and were not changed during it.

\subsection{Estimator Variants}

The four variants form a complexity ladder in which each rung adds exactly
one mechanism (Table~\ref{tab:estimators}). \textbf{T0} is a passthrough
control: one output per input, silent during silence, deliberately not
improved so that it reproduces the untreated failure modes. \textbf{T1} holds
the last valid estimate with exponential smoothing ($\alpha = 0.4$) for a
bounded hold duration, emits at a fixed rate, and expires into an explicit
empty output rather than a lingering track. \textbf{T2} is a linear
constant-velocity Kalman filter on the state $[c_x, c_y, \log w, \log h]$ and
its first derivatives, with measurement $z = [c_x, c_y, \log w, \log h]$,
timestamp-derived $\Delta t$, discrete white-acceleration process noise
($q = 0.30$), fixed measurement noise ($\sigma_c = 0.015$, $\sigma_s =
0.05$), and $\chi^2$ innovation gating at the 99\% level with four degrees of
freedom; rejected measurements do not stop prediction. Track expiry is
deterministic: prediction horizon, empty-message streak, exit from the image,
or maximum age. The logarithmic size parameterization keeps the scale
positive and makes the growth of apparent size approximately linear during a
straight approach.

\textbf{T3} is an adaptive-covariance framework in which the measurement
noise is modulated by detection features, $R = R_0 \exp(\theta^\top
\varphi)$, with $\varphi$ comprising inverse confidence, inverse square root
of apparent area, border proximity, and frame interval. \emph{We report T3 as
a framework, not as a result.} Its coefficients must be fitted to
synchronized detector-versus-reference residuals, which do not yet exist for
this platform; the shipped default is $\theta = 0$, under which T3 is exactly
T2, and a unit test asserts this identity. Consequently every T3 number in
this paper is a T2 number obtained through a second code path, and the
adaptive mechanism is neither validated nor claimed. The calibration
machinery itself was verified negatively: applied to synthetic homoscedastic
data it correctly recovers approximately zero feature dependence rather than
inventing structure.

%==============================================================================
% Temporal estimator ladder.  Source: docs/TEMPORAL_OBSTACLE_ESTIMATION.md,
% src/rov_obstacle_tracking.  Status: READY.
%==============================================================================
\begin{table}[t]
\centering
\caption{Temporal estimator variants. Each rung adds exactly one mechanism.
T3 ships with $\theta=0$ and is therefore \emph{identical} to T2 in every
experiment reported here; it is listed as implemented framework, not as an
evaluated method.}
\label{tab:estimators}
\footnotesize
\setlength{\tabcolsep}{3pt}
\renewcommand{\arraystretch}{1.2}
\begin{tabular}{@{}l p{20mm} p{17mm} p{22mm}@{}}
\toprule
& \textbf{Mechanism} & \textbf{Behavior in silence} & \textbf{Status} \\
\midrule
\textbf{T0} & Passthrough              & Silent            & Control baseline; deliberately not improved \\
\textbf{T1} & Hold + EMA ($\alpha=0.4$)& Hold, bounded 2.5\,s & Evaluated \\
\textbf{T2} & CV Kalman on $[c_x,c_y,\log w,\log h]$ + derivatives; $\chi^2$ gate & Predict, bounded 2.5\,s & Evaluated \\
\textbf{T3} & $R = R_0\exp(\theta^\top\varphi)$ over detection features & As T2 & \textbf{Framework only}: $\theta=0 \Rightarrow$ T3 $\equiv$ T2; calibration pending \\
\bottomrule
\end{tabular}
\end{table}


\subsection{Deterministic Replay and Controlled Disturbances}

Closed-loop runs alone cannot separate an estimator's behavior from the
planner's reaction to it, so the estimators are additionally evaluated on a
deterministic replay harness. Records carry the three upstream conditions as
first-class fields, and the harness interleaves message events with
fixed-rate ticks exactly as the online node does, so the same input file
yields bit-identical output. Reference values are consumed only by the
metrics layer, never by the estimator. Fifteen controlled cases D0--D14 span
clean tracking, silences of 0.25--2\,s, a silence beginning shortly after
maneuver commitment, a mid-maneuver silence, repeated silences, the
fresh-empty counterpart of the early silence, a single large outlier, a noise
burst, confidence collapse followed by dropout, timestamp jitter, a brief
false positive, and a legitimate disappearance. Metrics include overall and
prediction-only centre RMSE, availability, longest gap, reacquisition delay,
ghost-track time, normalized innovation squared, and empirical $1\sigma /
2\sigma$ coverage.

Hold duration and prediction horizon were both set to 2.5\,s using only the
development cases, whose longest dropout is 2.0\,s, and nothing was tuned on
the closed-loop evaluation scenarios. Dropouts longer than the horizon still
defeat T1 and T2 by construction, which is the intended behavior: the
estimators bridge gaps, they do not conjure obstacles.
